\chapter{PENUTUP}
\label{sec:chap5_tutup}

Bab ini menyajikan kesimpulan dan saran dari perancangan serta pengujian sistem \textit{Credit Risk Scoring} berbasis \textit{Machine Learning Operations} (MLOps).

\section{Kesimpulan}
\label{sec:sec5_kesimpulan}
Berdasarkan hasil penelitian dan pengujian, diperoleh beberapa kesimpulan sebagai berikut:

\begin{enumerate}
    \item Penelitian ini berhasil merancang pipeline MLOps untuk operasionalisasi model prediktif risiko kredit, mulai dari pembentukan data sintetis, pra-pemrosesan, pelatihan model, deployment API, interpretabilitas, hingga otomasi pelatihan ulang berkelanjutan (Continuous Training).

    \item Penggunaan data sintetis memungkinkan penelitian risiko kredit dilakukan tanpa menggunakan data nasabah riil. Dataset yang dibangun telah memuat atribut demografis, finansial, histori kredit, asuransi, produk pinjaman, serta parameter risiko yang relevan terhadap proses \textit{credit scoring}.

    \item Tahap pra-pemrosesan menggunakan WoE, IV, dan SMOTE membantu meningkatkan kualitas fitur serta menangani ketidakseimbangan kelas. Model XGBoost yang dikembangkan memperoleh AUC sebesar 0,8262, precision 0,8024, recall 0,7549, dan F1-score 0,7779.

    \item Integrasi SHAP, FastAPI, sistem grading, dan aturan bisnis membuat hasil prediksi lebih transparan dan operasional. Kategori Undefined juga membantu mencegah keputusan otomatis pada kasus yang membutuhkan pengecekan manual, seperti usia pelunasan yang melebihi batas kebijakan.

    \item Implementasi MLflow dan n8n menunjukkan bahwa proses pelacakan eksperimen, \textit{Continuous Training}, dan pemicuan retraining dapat diotomatisasi. Dengan demikian, sistem yang dibangun dapat menjadi prototipe awal pipeline MLOps untuk model risiko kredit.
\end{enumerate}

Secara keseluruhan, penelitian ini menunjukkan bahwa kombinasi model prediktif, interpretabilitas, aturan bisnis, dan otomasi MLOps dapat mendukung sistem penilaian risiko kredit yang lebih terstruktur. Namun, karena data yang digunakan masih bersifat sintetis, sistem ini tetap memerlukan validasi lebih lanjut sebelum diterapkan pada keputusan kredit riil.

\section{Saran}
\label{sec:sec5_saran}
Beberapa saran untuk pengembangan penelitian selanjutnya adalah sebagai berikut:

\begin{enumerate}
    \item Melakukan validasi menggunakan data historis riil yang telah dianonimkan agar performa model dapat diuji pada kondisi portofolio kredit yang sebenarnya.

    \item Mengembangkan proses pembentukan data sintetis dan membandingkan model XGBoost dengan pendekatan lain, seperti Logistic Regression, Random Forest, LightGBM, CatBoost, atau model berbasis \textit{scorecard}.

    \item Menyempurnakan aturan bisnis dan sistem grading dengan mempertimbangkan kebijakan internal bank, jenis produk kredit, batas usia, tenor, asuransi, serta skema \textit{co-borrower}.

    \item Memperkuat pipeline MLOps dengan audit trail, monitoring performa, validasi otomatis model baru, keamanan API, dan dashboard operasional agar sistem lebih siap digunakan pada lingkungan produksi.
\end{enumerate}
